\section*{ЗАКЛЮЧЕНИЕ}
\addcontentsline{toc}{section}{ЗАКЛЮЧЕНИЕ}

В ходе выполнения выпускной квалификационной работы была разработана и протестирована клиентская часть онлайн-площадки для аренды вещей, реализующая полный цикл взаимодействия пользователя с сервисом: от регистрации и просмотра каталога до подписания договора аренды и демонстрационной оплаты.
В процессе работы были решены все задачи, поставленные в техническом задании:

\begin{enumerate}
\item Проведён анализ предметной области онлайн-аренды и существующих решений. Рассмотрены классификация площадок по типу товаров, модели взаимодействия и типу арендодателей. Выявлены функциональные преимущества разрабатываемой системы по сравнению с аналогами Rentomania, Авито и Арендую.ру.
\item Разработана концептуальная модель клиентской части. Определены основные экраны, компоненты и сценарии взаимодействия пользователя с системой. Спроектированы варианты использования, охватывающий все ключевые операции сервиса.
\item Спроектирован пользовательский интерфейс, обеспечивающий взаимодействие с серверной частью посредством REST API. Определены все endpoint, используемые клиентской частью, и реализована централизованная функция apiRequest с автоматическим обновлением JWT-токена.
\item Разработана и протестирована клиентская часть веб-приложения. Реализованы модули авторизации с подтверждением email, каталога вещей с фильтрацией, бронирования с автоматическим расчётом стоимости с учётом наценки платформы и комиссии, встроенного чата с WebSocket, договоров аренды с демонстрационной ЭЦП, системы отзывов и уведомлений в реальном времени.
\end{enumerate}

Клиентская часть реализована в виде одностраничного приложения (SPA) на стандартных веб-технологиях HTML5, CSS3 и Vanilla JavaScript без использования фронтенд-фреймворков. Файл index.html содержит 7 626 строк кода, из которых 5 610 строк составляет программная логика JavaScript (207 функций). Для выбора дат аренды использована библиотека flatpickr, для хранения сессионных данных — localStorage.

Взаимодействие с серверной частью организовано через Django REST API для стандартных операций и через WebSocket-соединения (Django Channels, Redis) для чата и уведомлений в реальном времени. Реализован механизм автоматического переподключения WebSocket при разрыве связи и резервный опрос REST-endpoint уведомлений.

Проведено системное тестирование по 14 пользовательским сценариям в браузерах Google Chrome и Mozilla Firefox на настольных компьютерах и мобильных устройствах. Все сценарии дали результат, соответствующий ожидаемому.
